% GENERATED FILE. Edit canonical lessons, then run npm run generate:books.
% canonical-source-hash: fnv1a64:e432c3ef29db44f7
% canonical-lessons: RU-C148-zdorovyy, RU-C148-slabyy, RU-C148-silnyy, RU-C148-blednyy, RU-C148-ustalyy, RU-R148-first-pass-words-from-chapters-140-143, RU-R148-second-pass-words-from-chapters-144-148

\chapter{Healthy, Weak, Strong, Pale, Tired}
\label{ch:ru-healthy-weak-strong-pale-tired}
\hlchaptermodality{148}

\begin{chapteropening}
\textbf{By the end of this chapter:} \emph{I can say healthy, weak, strong, pale and tired.}

Complete the last lesson of chapter 148: I can say healthy, weak, strong, pale and tired.
\end{chapteropening}

\section[\texorpdfstring{\ru{здоровый} (zdoróvyy)}{здоровый (zdoróvyy)}]{\ru{здоровый} (zdoróvyy) — healthy}

\label{lesson:RU-C148-zdorovyy}



\begin{quote}
Before the new one: say the Russian for a runny nose, then the Russian for a pill.
\end{quote}

\subsection*{You'll want to know: \ru{здоровый}}
\textbf{\ru{здоровый}} (\emph{zdoróvyy}) — \textquotedblleft{}healthy\textquotedblright{}.

\begin{grammarlens}[title={What you've built}]
One more word for this chapter.
\end{grammarlens}

\subsection*{Your turn}
\begin{itemize}
\raggedright
  \item \emph{Say it:} \emph{zdoróvyy}
  \item \emph{Say it:} \emph{zdoróvyy}, once more
  \item \emph{Say it:} say \emph{zdoróvyy}
  \item \emph{Recall:} say the Russian for a runny nose, then the Russian for a pill, then say \emph{zdoróvyy} again
\end{itemize}

\begin{tcolorbox}[breakable,colback=teal!4,colframe=teal!35!black,title={Before you move on}]
What does \ru{здоровый} mean? (\textquotedblleft{}Healthy\textquotedblright{}.) Say it once more.
\end{tcolorbox}
\section[\texorpdfstring{\ru{слабый} (slábyy)}{слабый (slábyy)}]{\ru{слабый} (slábyy) — weak}

\label{lesson:RU-C148-slabyy}



\begin{quote}
Before the new one: say the Russian for a pill, then the Russian for healthy.
\end{quote}

\subsection*{You'll want to know: \ru{слабый}}
\textbf{\ru{слабый}} (\emph{slábyy}) — \textquotedblleft{}weak\textquotedblright{}.

\begin{grammarlens}[title={What you've built}]
One more word for this chapter.
\end{grammarlens}

\subsection*{Your turn}
\begin{itemize}
\raggedright
  \item \emph{Say it:} \emph{slábyy}
  \item \emph{Say it:} \emph{slábyy}, once more
  \item \emph{Say it:} say \emph{slábyy}
  \item \emph{Recall:} say the Russian for a pill, then the Russian for healthy, then say \emph{slábyy} again
\end{itemize}

\begin{tcolorbox}[breakable,colback=teal!4,colframe=teal!35!black,title={Before you move on}]
What does \ru{слабый} mean? (\textquotedblleft{}Weak\textquotedblright{}.) Say it once more.
\end{tcolorbox}
\section[\texorpdfstring{\ru{сильный} (sílnyy)}{сильный (sílnyy)}]{\ru{сильный} (sílnyy) — strong}

\label{lesson:RU-C148-silnyy}



\begin{quote}
Before the new one: say the Russian for healthy, then the Russian for weak.
\end{quote}

\subsection*{You'll want to know: \ru{сильный}}
\textbf{\ru{сильный}} (\emph{sílnyy}) — \textquotedblleft{}strong\textquotedblright{}.

\begin{grammarlens}[title={What you've built}]
One more word for this chapter.
\end{grammarlens}

\subsection*{Your turn}
\begin{itemize}
\raggedright
  \item \emph{Say it:} \emph{sílnyy}
  \item \emph{Say it:} \emph{sílnyy}, once more
  \item \emph{Say it:} say \emph{sílnyy}
  \item \emph{Recall:} say the Russian for healthy, then the Russian for weak, then say \emph{sílnyy} again
\end{itemize}

\begin{tcolorbox}[breakable,colback=teal!4,colframe=teal!35!black,title={Before you move on}]
What does \ru{сильный} mean? (\textquotedblleft{}Strong\textquotedblright{}.) Say it once more.
\end{tcolorbox}
\section[\texorpdfstring{\ru{бледный} (blédnyy)}{бледный (blédnyy)}]{\ru{бледный} (blédnyy) — pale}

\label{lesson:RU-C148-blednyy}



\begin{quote}
Before the new one: say the Russian for weak, then the Russian for strong.
\end{quote}

\subsection*{You'll want to know: \ru{бледный}}
\textbf{\ru{бледный}} (\emph{blédnyy}) — \textquotedblleft{}pale\textquotedblright{}.

\begin{grammarlens}[title={What you've built}]
One more word for this chapter.
\end{grammarlens}

\subsection*{Your turn}
\begin{itemize}
\raggedright
  \item \emph{Say it:} \emph{blédnyy}
  \item \emph{Say it:} \emph{blédnyy}, once more
  \item \emph{Say it:} say \emph{blédnyy}
  \item \emph{Recall:} say the Russian for weak, then the Russian for strong, then say \emph{blédnyy} again
\end{itemize}

\begin{tcolorbox}[breakable,colback=teal!4,colframe=teal!35!black,title={Before you move on}]
What does \ru{бледный} mean? (\textquotedblleft{}Pale\textquotedblright{}.) Say it once more.
\end{tcolorbox}
\section[\texorpdfstring{\ru{усталый} (ustályy)}{усталый (ustályy)}]{\ru{усталый} (ustályy) — tired}

\label{lesson:RU-C148-ustalyy}



\begin{quote}
Before the new one: say the Russian for strong, then the Russian for pale.
\end{quote}

\subsection*{You'll want to know: \ru{усталый}}
\textbf{\ru{усталый}} (\emph{ustályy}) — \textquotedblleft{}tired\textquotedblright{}.

\begin{grammarlens}[title={What you've built}]
One more word for this chapter.
\end{grammarlens}

\subsection*{Your turn}
\begin{itemize}
\raggedright
  \item \emph{Say it:} \emph{ustályy}
  \item \emph{Say it:} \emph{ustályy}, once more
  \item \emph{Say it:} say \emph{ustályy}
  \item \emph{Recall:} say the Russian for strong, then the Russian for pale, then say \emph{ustályy} again
\end{itemize}

\begin{tcolorbox}[breakable,colback=teal!4,colframe=teal!35!black,title={Before you move on}]
What does \ru{усталый} mean? (\textquotedblleft{}Tired\textquotedblright{}.) Say it once more.
\end{tcolorbox}
\section[(dialogue)]{Words from chapters 140-143 — a first pass}

\label{lesson:RU-R148-first-pass-words-from-chapters-140-143}



\begin{quote}
Say the words for pale and tired.
\end{quote}

\subsection*{Your turn}
Say these aloud:

\begin{itemize}
\raggedright
  \item to have lunch, to have breakfast, to have dinner, to do the laundry, to iron; to stroke
  \item to meet, to get acquainted, to invite, to thank, to apologise
  \item to carry; to wear, to hold, to throw, to catch, to push
  \item to fly, to ride, to walk, to come, to leave
\end{itemize}

\begin{tcolorbox}[breakable,colback=teal!4,colframe=teal!35!black,title={Before you move on}]
To have lunch? (\textbf{\ru{обедать}}.) To have breakfast? (\textbf{\ru{завтракать}}.) To have dinner? (\textbf{\ru{ужинать}}.) To do the laundry? (\textbf{\ru{стирать}}.) To iron; to stroke? (\textbf{\ru{гладить}}.) To meet? (\textbf{\ru{встречать}}.) To get acquainted? (\textbf{\ru{знакомиться}}.) To invite? (\textbf{\ru{приглашать}}.) To thank? (\textbf{\ru{благодарить}}.) To apologise? (\textbf{\ru{извиняться}}.) To carry; to wear? (\textbf{\ru{носить}}.) To hold? (\textbf{\ru{держать}}.) To throw? (\textbf{\ru{бросать}}.) To catch? (\textbf{\ru{ловить}}.) To push? (\textbf{\ru{толкать}}.) To fly? (\textbf{\ru{летать}}.) To ride? (\textbf{\ru{ездить}}.) To walk? (\textbf{\ru{ходить}}.) To come? (\textbf{\ru{приходить}}.) To leave? (\textbf{\ru{уходить}}.)
\end{tcolorbox}
\section[(dialogue)]{Words from chapters 144-148 — a second pass}

\label{lesson:RU-R148-second-pass-words-from-chapters-144-148}



\begin{quote}
Say the words for to be afraid and to please.
\end{quote}

\subsection*{Your turn}
Say these aloud:

\begin{itemize}
\raggedright
  \item to be afraid, to please, to hate, to praise, to scold
  \item to sign, to print, to keep; to save, to download, to be beaten
  \item to book, to turn, to cross, to stop, to park
  \item ill; a patient, a cold, flu, a runny nose, a pill
  \item healthy, weak, strong, pale, tired
\end{itemize}

\begin{tcolorbox}[breakable,colback=teal!4,colframe=teal!35!black,title={Before you move on}]
To be afraid? (\textbf{\ru{бояться}}.) To please? (\textbf{\ru{нравиться}}.) To hate? (\textbf{\ru{ненавидеть}}.) To praise? (\textbf{\ru{хвалить}}.) To scold? (\textbf{\ru{ругать}}.) To sign? (\textbf{\ru{подписывать}}.) To print? (\textbf{\ru{печатать}}.) To keep; to save? (\textbf{\ru{сохранять}}.) To download? (\textbf{\ru{скачивать}}.) To be beaten? (\textbf{\ru{проигрывать}}.) To book? (\textbf{\ru{бронировать}}.) To turn? (\textbf{\ru{поворачивать}}.) To cross? (\textbf{\ru{пересекать}}.) To stop? (\textbf{\ru{останавливаться}}.) To park? (\textbf{\ru{парковать}}.) Ill; a patient? (\textbf{\ru{больной}}.) A cold? (\textbf{\ru{простуда}}.) Flu? (\textbf{\ru{грипп}}.) A runny nose? (\textbf{\ru{насморк}}.) A pill? (\textbf{\ru{таблетка}}.) Healthy? (\textbf{\ru{здоровый}}.) Weak? (\textbf{\ru{слабый}}.) Strong? (\textbf{\ru{сильный}}.) Pale? (\textbf{\ru{бледный}}.) Tired? (\textbf{\ru{усталый}}.)
\end{tcolorbox}
